\section{Discussion and Conclusion}
\label{sec:conclusion}

We introduced \shadow, a $4{,}500$-image human-verified benchmark for geographic transfer in overhead shadow detection across three U.S.\ cities and two native resolutions, with a leave-one-city-out protocol and paired bootstrap reporting. Its three-probe framework indicates that the failure operates in feature space (Section~\ref{sec:feature-space}), appears as a class-asymmetric collapse of foreground confidence that $0/1$ metrics under-report by $2.3\times$ (Section~\ref{sec:sp-gap}), and is recoverable through an architecture-dependent component: decoder for DINOv3, encoder for OGLANet, both for MAMNet (Section~\ref{sec:architecture-locus}). Four of five selective-prediction findings replicate on INRIA building footprints. The gap remains significant in every $900$-image cell and on the external S-EO benchmark. As a case study, \sib{} has the smallest mean gap of the ten methods we evaluate ($6.42$ mIoU), level with FDA.

The benchmark covers three U.S.\ cities, RGB-only NAIP imagery, a single fall season, and four architectures. Non-grid morphologies, multispectral sensors, and untested architectures (e.g., SegFormer) remain open. We characterize where the failure occurs, how it distorts predictions, and which component recovers it, but do not isolate a causal mechanism that connects these characterizations. DINOv3's decoder locus also cannot be attributed to its ViT architecture alone, since no CNN in our tests matches its pretraining scale. We release \shadow, the LOCO splits, training code, and trained models as an evaluation instrument to be read with the three probes: its LOCO gaps measure geographic transfer risk in overhead dense prediction, and its two native resolutions support diagnosis of cross-resolution transfer failure. Together, the benchmark and framework give the field a shared instrument for the question that matters at deployment: not just whether a model works, but where and on which class it will fail when the city changes.

